% Studies-in-Algebra-and-Group-Theory - Lecture 1: Set Theory
% Author: S. D. V. Stephens
% Date: 2025-12-09
% Compile with: pdflatex -shell-escape

\documentclass{report}
\input{~/university/preamble.tex}
% preamble-darkmode.tex
% Dark mode extension for lecture notes
% Usage: % preamble-darkmode.tex
% Dark mode extension for lecture notes
% Usage: % preamble-darkmode.tex
% Dark mode extension for lecture notes
% Usage: \input{~/university/preamble-darkmode.tex} AFTER preamble.tex
%
% This provides:
% - Dark background with light text
% - Material Design color palette
% - Ocean blue gradient colors (p1-p9)
% - Enhanced TikZ/PGFPlots for dark backgrounds
% - qq tcolorbox environment
% - tkz-euclide for geometric constructions

% ===========================================
% DARK MODE PACKAGE
% ===========================================
\usepackage{darkmode}
\enabledarkmode

% ===========================================
% ADDITIONAL PACKAGES
% ===========================================
\usepackage{changepage}
\usepackage{ulem}
\usepackage{colortbl}
\usepackage{tkz-euclide}
\usepackage{capt-of}

% ===========================================
% EXTENDED TIKZ LIBRARIES
% ===========================================
\usetikzlibrary{patterns,3d,backgrounds}
\usepgfplotslibrary{groupplots}

% Upgrade pgfplots compatibility
\pgfplotsset{compat=1.18}

% ===========================================
% OCEAN BLUE GRADIENT (p1-p9)
% ===========================================
\definecolor{p1}{HTML}{caf0f8}
\definecolor{p2}{HTML}{ade8f4}
\definecolor{p3}{HTML}{90e0ef}
\definecolor{p4}{HTML}{48cae4}
\definecolor{p5}{HTML}{00b4d8}
\definecolor{p6}{HTML}{0096c7}
\definecolor{p7}{HTML}{0077b6}
\definecolor{p8}{HTML}{023e8a}
\definecolor{p9}{HTML}{03045e}

% Dark background color
\definecolor{pag}{HTML}{293133}

% ===========================================
% MATERIAL DESIGN COLORS
% ===========================================
% Note: myred, myyellow already defined in main preamble
% These are additional/alternate versions
\definecolor{matred}{HTML}{f44336}
\definecolor{matpink}{HTML}{e81e63}
\definecolor{matpurple}{HTML}{9c27b0}
\definecolor{matdeeppurple}{HTML}{673ab7}
\definecolor{matindigo}{HTML}{3f51b5}
\definecolor{matblue}{HTML}{2196f3}
\definecolor{matlightblue}{HTML}{03a9f4}
\definecolor{matcyan}{HTML}{00bcd4}
\definecolor{matteal}{HTML}{009688}
\definecolor{matgreen}{HTML}{4caf50}
\definecolor{matlightgreen}{HTML}{8bc34a}
\definecolor{matlime}{HTML}{cddc39}
\definecolor{matyellow}{HTML}{ffeb3b}
\definecolor{matamber}{HTML}{ffc107}
\definecolor{matorange}{HTML}{ff9800}
\definecolor{matdeeporange}{HTML}{ff5722}
\definecolor{matbrown}{HTML}{795548}
\definecolor{matgray}{HTML}{9e9e9e}
\definecolor{matbluegray}{HTML}{607d8b}

% ===========================================
% COLORLET FOR TIKZ
% ===========================================
\colorlet{xcol}{blue!60!black}

% ===========================================
% DARK MODE TCOLORBOX
% ===========================================
\newtcolorbox{qq}[2][]{%
    boxrule=0.75pt,
    sharp corners,
    colframe=white,
    colback=pag,
    coltext=white,
    #1,
}

% Dark definition box
\newtcolorbox{darkdfn}[2][]{%
    enhanced,
    boxrule=1pt,
    sharp corners,
    colframe=p5,
    colback=pag,
    coltext=white,
    coltitle=white,
    fonttitle=\bfseries,
    title=#2,
    #1,
}

% Dark theorem box
\newtcolorbox{darkthm}[2][]{%
    enhanced,
    boxrule=1pt,
    sharp corners,
    colframe=matpurple,
    colback=pag,
    coltext=white,
    coltitle=white,
    fonttitle=\bfseries,
    title=#2,
    #1,
}

% ===========================================
% TIKZ 3D FIX
% ===========================================
% Small fix for canvas is xy plane at z
% https://tex.stackexchange.com/a/48776/121799
\makeatletter
\tikzoption{canvas is xy plane at z}[]{%
    \def\tikz@plane@origin{\pgfpointxyz{0}{0}{#1}}%
    \def\tikz@plane@x{\pgfpointxyz{1}{0}{#1}}%
    \def\tikz@plane@y{\pgfpointxyz{0}{1}{#1}}%
    \tikz@canvas@is@plane}
\makeatother

% ===========================================
% CONDITIONAL LINE TIKZSET
% ===========================================
\tikzset{
    conditional line/.style 2 args={
        /utils/exec={\pgfmathsetmacro{\xcoordA}{#1}},
        /utils/exec={\pgfmathsetmacro{\xcoordB}{#2}},
        insert path={
            \ifdim\xcoordA pt>\xcoordB pt
            (\xcoordA,0) -- (\xcoordB,0)
            \fi
        },
    },
}

% ===========================================
% PGFPLOTS DARK MODE DEFAULTS
% ===========================================
\pgfplotsset{
    dark axis/.style={
        axis line style={white},
        tick style={white},
        label style={white},
        grid style={pag!70},
        every axis label/.append style={white},
        every tick label/.append style={white},
    },
}

% ===========================================
% TIKZ EXTERNALIZATION (for fast compilation)
% ===========================================
% This creates .md5 cache files for figures
% Requires: pdflatex -shell-escape
%
% To enable, add AFTER loading this file:
%   \enabletikzexternalize
%
\newcommand{\enabletikzexternalize}{%
  \usetikzlibrary{external}%
  \tikzexternalize[prefix=figures/]%
}

% ===========================================
% CONVENIENT SHORTCUTS FOR DARK PLOTS
% ===========================================
\newcommand{\darkplot}[1]{%
    \begin{tikzpicture}
    \begin{axis}[
        dark axis,
        grid=both,
        major grid style={line width=.2pt,draw=pag!70},
        #1
    ]
}


% ===========================================
% QUICK GEOMETRY MACROS (tkz-euclide)
% ===========================================
% Draw a point: \point{name}{x}{y}
\newcommand{\gpoint}[3]{\tkzDefPoint(#2,#3){#1}}

% Draw line through points: \gline{A}{B}
\newcommand{\gline}[2]{\tkzDrawLine[color=p4](#1,#2)}

% Draw circle: \gcircle{center}{radius}
\newcommand{\gcircle}[2]{\tkzDrawCircle[color=p5](#1,#2)}

% Mark angle: \gangle{A}{B}{C}
\newcommand{\gangle}[3]{\tkzMarkAngle[color=p3,size=0.5](#1,#2,#3)}

% ===========================================
% USAGE EXAMPLE
% ===========================================
% In your document:
%
% \begin{tikzpicture}
%   \begin{axis}[dark axis, ...]
%     \addplot[p4, thick] {sin(deg(x))};
%   \end{axis}
% \end{tikzpicture}
%
% Or with qq box:
% \begin{qq}{My Title}
%   Content here in dark box
% \end{qq}
 AFTER preamble.tex
%
% This provides:
% - Dark background with light text
% - Material Design color palette
% - Ocean blue gradient colors (p1-p9)
% - Enhanced TikZ/PGFPlots for dark backgrounds
% - qq tcolorbox environment
% - tkz-euclide for geometric constructions

% ===========================================
% DARK MODE PACKAGE
% ===========================================
\usepackage{darkmode}
\enabledarkmode

% ===========================================
% ADDITIONAL PACKAGES
% ===========================================
\usepackage{changepage}
\usepackage{ulem}
\usepackage{colortbl}
\usepackage{tkz-euclide}
\usepackage{capt-of}

% ===========================================
% EXTENDED TIKZ LIBRARIES
% ===========================================
\usetikzlibrary{patterns,3d,backgrounds}
\usepgfplotslibrary{groupplots}

% Upgrade pgfplots compatibility
\pgfplotsset{compat=1.18}

% ===========================================
% OCEAN BLUE GRADIENT (p1-p9)
% ===========================================
\definecolor{p1}{HTML}{caf0f8}
\definecolor{p2}{HTML}{ade8f4}
\definecolor{p3}{HTML}{90e0ef}
\definecolor{p4}{HTML}{48cae4}
\definecolor{p5}{HTML}{00b4d8}
\definecolor{p6}{HTML}{0096c7}
\definecolor{p7}{HTML}{0077b6}
\definecolor{p8}{HTML}{023e8a}
\definecolor{p9}{HTML}{03045e}

% Dark background color
\definecolor{pag}{HTML}{293133}

% ===========================================
% MATERIAL DESIGN COLORS
% ===========================================
% Note: myred, myyellow already defined in main preamble
% These are additional/alternate versions
\definecolor{matred}{HTML}{f44336}
\definecolor{matpink}{HTML}{e81e63}
\definecolor{matpurple}{HTML}{9c27b0}
\definecolor{matdeeppurple}{HTML}{673ab7}
\definecolor{matindigo}{HTML}{3f51b5}
\definecolor{matblue}{HTML}{2196f3}
\definecolor{matlightblue}{HTML}{03a9f4}
\definecolor{matcyan}{HTML}{00bcd4}
\definecolor{matteal}{HTML}{009688}
\definecolor{matgreen}{HTML}{4caf50}
\definecolor{matlightgreen}{HTML}{8bc34a}
\definecolor{matlime}{HTML}{cddc39}
\definecolor{matyellow}{HTML}{ffeb3b}
\definecolor{matamber}{HTML}{ffc107}
\definecolor{matorange}{HTML}{ff9800}
\definecolor{matdeeporange}{HTML}{ff5722}
\definecolor{matbrown}{HTML}{795548}
\definecolor{matgray}{HTML}{9e9e9e}
\definecolor{matbluegray}{HTML}{607d8b}

% ===========================================
% COLORLET FOR TIKZ
% ===========================================
\colorlet{xcol}{blue!60!black}

% ===========================================
% DARK MODE TCOLORBOX
% ===========================================
\newtcolorbox{qq}[2][]{%
    boxrule=0.75pt,
    sharp corners,
    colframe=white,
    colback=pag,
    coltext=white,
    #1,
}

% Dark definition box
\newtcolorbox{darkdfn}[2][]{%
    enhanced,
    boxrule=1pt,
    sharp corners,
    colframe=p5,
    colback=pag,
    coltext=white,
    coltitle=white,
    fonttitle=\bfseries,
    title=#2,
    #1,
}

% Dark theorem box
\newtcolorbox{darkthm}[2][]{%
    enhanced,
    boxrule=1pt,
    sharp corners,
    colframe=matpurple,
    colback=pag,
    coltext=white,
    coltitle=white,
    fonttitle=\bfseries,
    title=#2,
    #1,
}

% ===========================================
% TIKZ 3D FIX
% ===========================================
% Small fix for canvas is xy plane at z
% https://tex.stackexchange.com/a/48776/121799
\makeatletter
\tikzoption{canvas is xy plane at z}[]{%
    \def\tikz@plane@origin{\pgfpointxyz{0}{0}{#1}}%
    \def\tikz@plane@x{\pgfpointxyz{1}{0}{#1}}%
    \def\tikz@plane@y{\pgfpointxyz{0}{1}{#1}}%
    \tikz@canvas@is@plane}
\makeatother

% ===========================================
% CONDITIONAL LINE TIKZSET
% ===========================================
\tikzset{
    conditional line/.style 2 args={
        /utils/exec={\pgfmathsetmacro{\xcoordA}{#1}},
        /utils/exec={\pgfmathsetmacro{\xcoordB}{#2}},
        insert path={
            \ifdim\xcoordA pt>\xcoordB pt
            (\xcoordA,0) -- (\xcoordB,0)
            \fi
        },
    },
}

% ===========================================
% PGFPLOTS DARK MODE DEFAULTS
% ===========================================
\pgfplotsset{
    dark axis/.style={
        axis line style={white},
        tick style={white},
        label style={white},
        grid style={pag!70},
        every axis label/.append style={white},
        every tick label/.append style={white},
    },
}

% ===========================================
% TIKZ EXTERNALIZATION (for fast compilation)
% ===========================================
% This creates .md5 cache files for figures
% Requires: pdflatex -shell-escape
%
% To enable, add AFTER loading this file:
%   \enabletikzexternalize
%
\newcommand{\enabletikzexternalize}{%
  \usetikzlibrary{external}%
  \tikzexternalize[prefix=figures/]%
}

% ===========================================
% CONVENIENT SHORTCUTS FOR DARK PLOTS
% ===========================================
\newcommand{\darkplot}[1]{%
    \begin{tikzpicture}
    \begin{axis}[
        dark axis,
        grid=both,
        major grid style={line width=.2pt,draw=pag!70},
        #1
    ]
}


% ===========================================
% QUICK GEOMETRY MACROS (tkz-euclide)
% ===========================================
% Draw a point: \point{name}{x}{y}
\newcommand{\gpoint}[3]{\tkzDefPoint(#2,#3){#1}}

% Draw line through points: \gline{A}{B}
\newcommand{\gline}[2]{\tkzDrawLine[color=p4](#1,#2)}

% Draw circle: \gcircle{center}{radius}
\newcommand{\gcircle}[2]{\tkzDrawCircle[color=p5](#1,#2)}

% Mark angle: \gangle{A}{B}{C}
\newcommand{\gangle}[3]{\tkzMarkAngle[color=p3,size=0.5](#1,#2,#3)}

% ===========================================
% USAGE EXAMPLE
% ===========================================
% In your document:
%
% \begin{tikzpicture}
%   \begin{axis}[dark axis, ...]
%     \addplot[p4, thick] {sin(deg(x))};
%   \end{axis}
% \end{tikzpicture}
%
% Or with qq box:
% \begin{qq}{My Title}
%   Content here in dark box
% \end{qq}
 AFTER preamble.tex
%
% This provides:
% - Dark background with light text
% - Material Design color palette
% - Ocean blue gradient colors (p1-p9)
% - Enhanced TikZ/PGFPlots for dark backgrounds
% - qq tcolorbox environment
% - tkz-euclide for geometric constructions

% ===========================================
% DARK MODE PACKAGE
% ===========================================
\usepackage{darkmode}
\enabledarkmode

% ===========================================
% ADDITIONAL PACKAGES
% ===========================================
\usepackage{changepage}
\usepackage{ulem}
\usepackage{colortbl}
\usepackage{tkz-euclide}
\usepackage{capt-of}

% ===========================================
% EXTENDED TIKZ LIBRARIES
% ===========================================
\usetikzlibrary{patterns,3d,backgrounds}
\usepgfplotslibrary{groupplots}

% Upgrade pgfplots compatibility
\pgfplotsset{compat=1.18}

% ===========================================
% OCEAN BLUE GRADIENT (p1-p9)
% ===========================================
\definecolor{p1}{HTML}{caf0f8}
\definecolor{p2}{HTML}{ade8f4}
\definecolor{p3}{HTML}{90e0ef}
\definecolor{p4}{HTML}{48cae4}
\definecolor{p5}{HTML}{00b4d8}
\definecolor{p6}{HTML}{0096c7}
\definecolor{p7}{HTML}{0077b6}
\definecolor{p8}{HTML}{023e8a}
\definecolor{p9}{HTML}{03045e}

% Dark background color
\definecolor{pag}{HTML}{293133}

% ===========================================
% MATERIAL DESIGN COLORS
% ===========================================
% Note: myred, myyellow already defined in main preamble
% These are additional/alternate versions
\definecolor{matred}{HTML}{f44336}
\definecolor{matpink}{HTML}{e81e63}
\definecolor{matpurple}{HTML}{9c27b0}
\definecolor{matdeeppurple}{HTML}{673ab7}
\definecolor{matindigo}{HTML}{3f51b5}
\definecolor{matblue}{HTML}{2196f3}
\definecolor{matlightblue}{HTML}{03a9f4}
\definecolor{matcyan}{HTML}{00bcd4}
\definecolor{matteal}{HTML}{009688}
\definecolor{matgreen}{HTML}{4caf50}
\definecolor{matlightgreen}{HTML}{8bc34a}
\definecolor{matlime}{HTML}{cddc39}
\definecolor{matyellow}{HTML}{ffeb3b}
\definecolor{matamber}{HTML}{ffc107}
\definecolor{matorange}{HTML}{ff9800}
\definecolor{matdeeporange}{HTML}{ff5722}
\definecolor{matbrown}{HTML}{795548}
\definecolor{matgray}{HTML}{9e9e9e}
\definecolor{matbluegray}{HTML}{607d8b}

% ===========================================
% COLORLET FOR TIKZ
% ===========================================
\colorlet{xcol}{blue!60!black}

% ===========================================
% DARK MODE TCOLORBOX
% ===========================================
\newtcolorbox{qq}[2][]{%
    boxrule=0.75pt,
    sharp corners,
    colframe=white,
    colback=pag,
    coltext=white,
    #1,
}

% Dark definition box
\newtcolorbox{darkdfn}[2][]{%
    enhanced,
    boxrule=1pt,
    sharp corners,
    colframe=p5,
    colback=pag,
    coltext=white,
    coltitle=white,
    fonttitle=\bfseries,
    title=#2,
    #1,
}

% Dark theorem box
\newtcolorbox{darkthm}[2][]{%
    enhanced,
    boxrule=1pt,
    sharp corners,
    colframe=matpurple,
    colback=pag,
    coltext=white,
    coltitle=white,
    fonttitle=\bfseries,
    title=#2,
    #1,
}

% ===========================================
% TIKZ 3D FIX
% ===========================================
% Small fix for canvas is xy plane at z
% https://tex.stackexchange.com/a/48776/121799
\makeatletter
\tikzoption{canvas is xy plane at z}[]{%
    \def\tikz@plane@origin{\pgfpointxyz{0}{0}{#1}}%
    \def\tikz@plane@x{\pgfpointxyz{1}{0}{#1}}%
    \def\tikz@plane@y{\pgfpointxyz{0}{1}{#1}}%
    \tikz@canvas@is@plane}
\makeatother

% ===========================================
% CONDITIONAL LINE TIKZSET
% ===========================================
\tikzset{
    conditional line/.style 2 args={
        /utils/exec={\pgfmathsetmacro{\xcoordA}{#1}},
        /utils/exec={\pgfmathsetmacro{\xcoordB}{#2}},
        insert path={
            \ifdim\xcoordA pt>\xcoordB pt
            (\xcoordA,0) -- (\xcoordB,0)
            \fi
        },
    },
}

% ===========================================
% PGFPLOTS DARK MODE DEFAULTS
% ===========================================
\pgfplotsset{
    dark axis/.style={
        axis line style={white},
        tick style={white},
        label style={white},
        grid style={pag!70},
        every axis label/.append style={white},
        every tick label/.append style={white},
    },
}

% ===========================================
% TIKZ EXTERNALIZATION (for fast compilation)
% ===========================================
% This creates .md5 cache files for figures
% Requires: pdflatex -shell-escape
%
% To enable, add AFTER loading this file:
%   \enabletikzexternalize
%
\newcommand{\enabletikzexternalize}{%
  \usetikzlibrary{external}%
  \tikzexternalize[prefix=figures/]%
}

% ===========================================
% CONVENIENT SHORTCUTS FOR DARK PLOTS
% ===========================================
\newcommand{\darkplot}[1]{%
    \begin{tikzpicture}
    \begin{axis}[
        dark axis,
        grid=both,
        major grid style={line width=.2pt,draw=pag!70},
        #1
    ]
}


% ===========================================
% QUICK GEOMETRY MACROS (tkz-euclide)
% ===========================================
% Draw a point: \point{name}{x}{y}
\newcommand{\gpoint}[3]{\tkzDefPoint(#2,#3){#1}}

% Draw line through points: \gline{A}{B}
\newcommand{\gline}[2]{\tkzDrawLine[color=p4](#1,#2)}

% Draw circle: \gcircle{center}{radius}
\newcommand{\gcircle}[2]{\tkzDrawCircle[color=p5](#1,#2)}

% Mark angle: \gangle{A}{B}{C}
\newcommand{\gangle}[3]{\tkzMarkAngle[color=p3,size=0.5](#1,#2,#3)}

% ===========================================
% USAGE EXAMPLE
% ===========================================
% In your document:
%
% \begin{tikzpicture}
%   \begin{axis}[dark axis, ...]
%     \addplot[p4, thick] {sin(deg(x))};
%   \end{axis}
% \end{tikzpicture}
%
% Or with qq box:
% \begin{qq}{My Title}
%   Content here in dark box
% \end{qq}


% TikZ Externalization (for fast recompilation)

\course{Studies-in-Algebra-and-Group-Theory}

\begin{document}

\lecture{1}{Set Theory}

\section{Ordered Pairs}
We begin with a discussion of ordered pairs. Consider a set of numbers \(\{1,2,3,4\}\). For any set there is no particular order that triumphs over another, i.e., \(\{1,2,3,4\} = \{1,3,2,4\} = \{3,4,2,1\} =\ldots\), an obvious truth. Say we want \(\{1,2,3,4\}\) to be in the order of \(\{3,2,1,4\}\) for some arbitrary reason; we may then consider the collection,
\[
  \mfC = \{\{3\} , \{3,2\} ,\{3,2,1\} , \{3,2,1,4\}     \}
.\]
It becomes obvious what is happening here: the method by which we get order is dependent on the number of iterations in which a particular element appears in any number of sets in a given collection (in comparison to the other elements in the other sets of the respective collection).

Then consider this process but for a nested collection (a collection within a collection) which contains two sets, namely the arbitrary elements \(a,b\) in opposing orders, and notice that it gives us the definition of a 2-tuple (ordered pair),
\[
  \mfC_{1} = \{\{a\} ,\{a,b\}   \} \coloneqq (a,b) \, \text{ and }\, \mfC_{2}  = \{\{b\} ,\{b,a\}   \} \coloneqq  (b,a)
.\]
\nt{This form of generating or representing ordered pairs is called the \emph{Kuratowski} of ordered pairs.
}
This, of course, naturally extends to greater \(n\)-tuples. There are some obvious consequences of this that will not be gone over here.

The \textbf{Cartesian product} is a set of ordered pairs which can be given by taking power sets.

\dfn{Cartesian Product}{The Cartesian product is the cross product of two sets given by
\[ A \times  B = \{(a,b)\mid a\in A, b\in B\}  \}   ,\] where, by the Kuratowski encoding,
\[  (a,b)\coloneqq \{\{a\},\{a,b\}\} \in  \mcP(\mcP(A \cup B))  .\]
}
\section{Relations}
The review or ordered pairs leads us to our next discussion on relations. A relation in more colloquial terms is a dynamic or characteristic of a given element shared with another given element such that they are arranged in the form of ordered pairs. Think of equality, then think of equality as equating \((2+3,1+4)\), or perhaps something more interesting, the set of all \((x,y)\) such that \(x\) is a man, \(y\) is an enby (slay), and \(x\) is married to \(y\) (many such cases). We can then think of relations as a set of ordered pairs \(R \subset A\times B\). Consider the following examples:

\begin{enumerate}
\item \(R \coloneqq \{(x,y)\in \mathbb{N}^2 \mid y-x>0\} \iff xRy  \equiv y>x\), the relation of greater than;
\item \(R\coloneqq \{ (a,B) \in A \times \mathcal{P}(A) \mid a \in B\} \iff aRB \equiv a \in B\), the relation of \emph{belonging} (being an element of);
\item \(R\coloneqq \{(c,d)\in \{\text{all cats}\} \times  \{\text{all dogs}\}  \mid c,d \text{ live in the same house}\}\).
\end{enumerate}

Naturally, there are an infinite number of possible relations, however there is a particular class of these which we are interested in for the time being---that which relates two equivalent elements (the definition of what equivalence means is, for now, notably ambiguous).

If a relation \(R\) exists s.t. \(xRx\) for any \(x \in  X\), then we say \(R\) is reflexive. If \(xRy \iff yRx\), it is symmetric, and if \(xRy, \,  yRz \implies xRz\), it is transitive. If a given relation \(R\) satisfies all three of these properties (reflexivity, symmetry, transitivity), then we say it is an \textbf{\emph{equivalence relation}}.

\dfn{Equivalence relations}{An equivalence relation, usually denoted \(\sim\), is a relation that is reflexive, symmetric, and transitive \(\forall (a,b) \in  A \times  B\).
}

\qs{}{Consider each of the three qualifiers (axioms) for the existence of an equivalence relation. Find an explicit relation that has one property but not the other for each of the three axioms.
}

\begin{solution}
I apologize for the terse abstraction involved in this first example, but let
\[ R\coloneqq \{(1,1),(2,2),(3,3)(1,2),(2,3)\in \{1,2,3\}^2  \}  ,\]
notice that \((1,1),(2,2),(3,3)\) is true for \(R\), thus reflexivity holds; however, it is quick to see that \((1,2)\) exists in \(R\) but \((2,1)\not\in R\), thus symmetry cannot hold. For transitivity, notice that we have \((1,2) \text{ and }  (2,3)\), so we can set up the first step, but \((1,3)\not\in  R\)! Thus transitivity fails. This example is a reminder that we need not think of concrete properties that can be applied in front of our eyes (shapes, dogs, even equality!) in order to have or manipulate some sort of mathematical relation. Next, consider a the relation \(R\coloneqq \{(a,b)\mid a \neq  b\}\). First, \(aRa\) is categorically untrue thus reflexivity means nothing to us. Symmetry, however, tells us that if \(aRb\) then \(bRa\), which is true! Transitivity may seems sneakily true at first, but notice that \(aRb, \, bRa \implies aRa\) is not true (we simply replaced \(c\) with another instance of \(a\)). Finally, transitivity as a singular property is pretty straightforward. Consider \(x>y\), notice that \(x>x \implies x<x\) is categorically untrue since \(x\) can never be greater or less than itself, additionally, \(x>y\) never implies \(y>x\). For transitivity however, notice that \(x>y, \, y>z \implies x>z\) is true!
\end{solution}


Returning to our discussion of relations and partitions we can have this thing called an \emph{induced relation}. Let \(R\coloneqq X / \mfC\), we say that if \(x,y \in  \mfC\) and \(xX / \mfC y\), then \(X / \mfC\) is the relation \emph{induced} by the partition \(\mfC\).

\mprop{Relationship between equivalence relations, classes, and set partitions}{Explicitly the relationship between the two is given as follows: if \(R\) is an equivalence relation on \(X\), then the set of equivalence classes form a partition of \(X\) that induces the relation \(R\), and if \(\mfC\) is a partition of \(X\), then the induced relation is an equivalence relation whose set of equivalence classes is exactly \(\mfC\).
}
\section{Functions}

First I want to remind us of some notation and terminology things (mathematicians are \textbf{\emph{very}} specific):
\nt{To say a function is \emph{on} \(X\) \emph{into} \(Y\), means that the function is from \(X\) to \(Y\)
}

\dfn{Functions}{\{We all know what functions are at some level, but here were are going to be extraordinarily explicit (though you likely have already heard it in this form, regardless). A function on \(X\) into \(y\) is a relation (\(R\)) given by a letter like \(f,g,h,\varphi,\ldots\), such that:
\begin{enumerate}
\item \(\dom f = X\);
\item if \((x,y)\in f\) and \((x,z) \in f\), then \(y=z\) (uniqueness of \(y\) for each \(x\)); additionally, \(f(x)=y\) is conventional notation (instead of \((x,y) \in f\) or \(xfy\));
\item \(y\) is called the \textbf{\emph{value}} that \(f\) \textbf{\emph{assumes}} at the \textbf{\emph{argument}} \(x\); i.e., \(f\) \textbf{\emph{maps} sends/ transforms/} \(x\) to/ into \(y\). We often use the terms map, mapping, transformation, correspondence, operator, and function interchangeably.
\end{enumerate}
We denote this correspondence via
\[ f : X \to  Y    .\]
}

Finally, we note that the set of all functions from \(X\) to \(Y\) is a subset of the power set \(\mcP (X \times Y)\) and is denoted \(Y^X\).

Something to note is that a function is not an active entity despite action-invoking words such as \emph{transforms}; it merely \emph{is}. The use of the word function to describe the undefined object that is somehow active in relation to a graph is nevertheless a static set, similar to that of a directory of names that correspond to addresses.
Reminder that while the \(\dom f = X\), it need not be the case that \(\ran f =Y\), instead the range is reserved for the subset of all \(y \in Y\) that have a corresponding \(x\) such that \(f(x)=y\), i.e., is mapped to by some \(x\). We call this property being \textbf{\emph{onto}} and we say \(f\) maps \(X\) \textbf{\emph{onto}} Y; more commonly however we say that \(f\) is \textbf{\emph{surjective}}---we explicitly write this as \(\forall y \in  Y, \exists  x \in  X : f(x)=y\)---and we denote this surjectivity via
\[ f : X \surjto Y,\]
and if \(A\subset X\), we may consider all \(y\in  Y\) such that there is an \(x\in A\) such that \(f(x_A)=y\) (here I use \(x_A\) as crude shorthand for \(f(x)\) s.t. \(x \in A\), where Halmos uses the more primitive \(f(A)\)), the set of those \(y \in Y\) is known as the \textbf{\emph{image}} of \(f\) under \(A\). If it is the case that \(X \subset  Y\), the function \(f\) defined by \(f(x)=x , \, \forall x \in X\), we call this the \textbf{\emph{inclusion map} embedding/ injection/} or we say \(f\) is \textbf{\emph{one-to-one}}; we explicitly write this as \(\forall a\neq b, \, f(a)\neq f(b)\). We denote injectivity via
\[ f : A \injto Y.\]
A special case of injectivity comes when we inject \(X \injto X\), relation-wise this is equality in \(X\), language-wise, we call this the \textbf{\emph{identity map}} on \(X\).

\mprop{Restriction and extensions}{Consider a function \(f:Y\to Z\) and a set \(X  \subset  Y\). We say that there is a \textbf{\emph{natural way}} of constructing a new function \(g : X \to  Z\); let \(g(x)=f(x), \forall x \in  X\), we then say \(g\) is a \textbf{\emph{restriction}} of \(f\) to \(X\) and \(f\) is the \textbf{\emph{extension}} of \(g\) to \(Y\). We denote this as \(g=f|X\) where we can express the definition as \((f|X)(x)=f(x) \forall x \in X\) and \(\ran(f|X)=f(x_X)\).
}

There are a few other terminologies that follow from all that has been mentioned so far, and I will go through them rather quickly:

\dfn{Projection}{\(f\) (also \(\pr\) is a \textbf{\emph{projection}} from \(X \times Y\) onto \(X\) if if \(f(x,y)=x\); similarly, if \(R=X \times  Y\), then instead of being a projection of \(R\) onto \(X\), it is now the range of the projection \(f\) (\(\ran (\pr )\)).
}

\dfn{Bijection}{Expounding on our notion of inclusion mappings and surjections from earlier, \textbf{\emph{bijectivity}}---one-to-one correspondence-ness---is a function that is both injective and surjective.
}
Bijections are derived from an equivalence relation \(R\) in \(X\) with \(aRb \iff f(a)=f(b)\) for \(a,b \in X\); so given a function/ set \(g(y)=\{x \in X \mid f(x)=y, \, \forall  y \in  Y\}\), the equivalence relation \(R\) tells us that \(g(y)\) is an equivalence class of the relation \(R\); that is to say, we define the function as \(g: Y \surjto X / R\), notice that for \(g\), if \(c\neq d, \, c,d \in Y\), due to the nature of equivalence classes (remember, all the classes are disjoint), \(g(c)\neq g(d)\); notice this is the exact explicit definition of injectivity, thus this is a surjective-injective, i.e., bijective mapping. Although there is no universially agreed on notation for bijection, we shall denote this mapping as \(g : Y \leftrightarrow X / R\).

\dfn{Characteristic functions}{If \(A \subset X\) the characteristic function of \(A\) is the function\(\mcX : X \to 2\) given by
 \[ \mcX(x)=\begin{cases}
      1 & \text{if } x \in  A \\
      0 & \text{if } x\in X-A
  \end{cases} \],
  however, we may write \(\mcX_A\) to indicate that it depends on \(A\). The function that assigns to each \(A \subset X\), i.e., \(\mcP(X)\), the characteristic function of \(A  \in  2^X\) is a bijection, i.e., \(f : \mcP(X) \leftrightarrow 2^X\).
}
\pagebreak
\qs{}{Concerning bijectivity (point iii.), the examples have all the inclusion maps as one-to-one, however, except in a few trivial special cases, the projections are not. What are these special cases?
}
\begin{solution}
Consider a function's projection as the cannonical map \(\gamma  : X \to X / R\), the following cases are trivial examples of non-injective projection maps:
\begin{enumerate}
\item \(X = \emptyset\), vacuously fulfilled (it matters not what \(R\) is);
\item \(R \coloneqq \left\{ (x,x) \mid x \in  X \right\}\), so every equivalence class is a singleton.
\end{enumerate}
\end{solution}

\qs{}{Prove:
\begin{enumerate}
\item \(Y^{\emptyset}\) has exactly one element, namely \(\emptyset\), regardless of if \(Y\) is empty or not.
\item If \(X\) is non-empty, then \(\emptyset^X\) is empty.
\end{enumerate}
}


\begin{solution}
\begin{enumerate}
\item Given that \(Y^{\emptyset} \subset \mcP (\emptyset \times Y)\), and the definition of Cartesian sets has \textbf{and} as a qualifier, i.e., \(A \times  B = \{(a,b) \mid a \in A \text{\textbf{ AND }} b \in  B \}\), we can see that \(\emptyset \times  Y\) contains in fact no elements which satisfy this, and thus, it is the empty set. So \(Y^{\emptyset}  \subset  \mcP(\emptyset) = \{\emptyset\}\), thus the only element that can be in \(Y^{\emptyset}\) is the singleton \(\{\emptyset\}\). This is a vacuous satisfaction and is non-intuitive; some intuition, perhaps, may be gained but thinking ``\emph{the only thing nothing can map to is nothing},'' or by taking the cardinality of these two sets with each other given by \(|\emptyset|\cdot |Y|=0|Y|=0\) (though this is mainly for intuition).
\item \(\emptyset^X\) means that something must map to the emptyset, meaning there must be an element within the emptyset, but this cannot be true! Thus, by definition, the set must be empty.\}
\end{enumerate}
\end{solution}
\section{Families}

Forget literally everything you know about notational convention! Why? Idk, that's just how mathematicians decided to approach this topic---the family (\emph{awwwwww}). Consider a function \(x\) (yes, this is strange) s.t. \(x : I \to  X\). We call \(I\) the \textbf{\emph{index set}}, and thereby each \(i \in  I\) is an \textbf{\emph{index}}; the range of the function is then called the \textbf{\emph{indexed set}} (why, this is so stupid?!); the function has a special name, being called a \textbf{\emph{family}}, and the value of \(x\) at a given \(i\) is called a \textbf{\emph{term}} of the family (denoted \(x_{i}\)). Other even more bewildering notation may be used such as \(\{x_{i}\}\) being a family in \(X\) or, more horrifically, a family \(\{A_i\}\) of subsets of \(X\), being a function \(A : I \to  \mcP(X)\) (now that I think of it, this notation, sickeningly, is becoming somewhat preferable).

Given that \(\ran (\{A_i\}= \bigcup_{i \in  I} A_i  )\), we notice that every arbitrary collection of sets is in fact the range of some arbitrary family: consider \(I\) and \(X\) equal to \(\mfC\) s.t. \(\{A_i\} : \mfC \to  \mfC\). I am realizing it is the terminology which is so obtuse, that what I do not like.

We may think of \(\bigcup_{i \in  I}  A_i\) as \(\bigcup^n_{i = 1}  A_i\) (similar to what we have seen for summations and products), where \(n\) is the order of \(i\) and we assume \(1\) is the first element in \(i\).

\ex{}{For some index set \(i=\{1,2,3\}\), and an indexed set \(A_i\) we have \(\bigcup^3_{i = 1}  = A_{1} \cup A_{2} \cup A_{3}\). Another example may be given as follows: let \(I = \mathbb{N}\), and \(A_i = \{-i,0,i\}\), that is to say for every \(i\) we have \(A_{1}=\{-1,0,1\}, A_{2}=\{-2,0,2\},\ldots\), and so on. Notice that the infinite union of \(A_i\) produces \(\mathbb{Z}\) and the infinite intersection produces \(\{0\}\).
}

\ex{}{Let us now consider an uncountably infinite indexing set it a plane \(I = [-1,1], A_i=\{i\}\times [0,1] \subset  \mathbb{R}^2\). Notice this is going to be a subset of ordered pairs, i.e., \(\{(i,y) \mid  0 \le y \le 1\}\). Geometrically, the union over all \(i\in I\) of \(A_i\) gives us a rectangle in \(\mathbb{R}^2\) defined by \([-1,1]\times [0,1]\); similarly the intersection of these sets, notice, are intersections of finite parallel lines, thus there exists no element in common with these sets, thus it is \(\emptyset\).
}

\thm{Associative law of unions}{Let \(\{I_j\} : J \to  K \mid  K = \bigcup_{j\in  I}\) and \(\{A_k\}: K \to X\), then
\[
  \bigcup_{k \in  K}A_k=\bigcup_{j \in  J}\left( \bigcup_{i \in I_j} A_i  \right)
.\]
}

\qs{}{Prove the theorem above.
}

Sorry! I am leaving this one as an exercise for the reader (sorrryyyyyy:) )


De Morgan's laws apply as standard for index/ indexed sets and their unions. From such, it is easy to see that if \(\{A_i\}, \{B_i\}\) are families of sets, then \(\bigcup_{i} A_i \cap \bigcup_j B_j = \bigcup_{i,j} (A_i \cap  B_j)\), a similarly for \(\bigcap_{i,j}  (A_i  \cup B_j)\).

Let us consider a special type of family: let \(I=\{a,b\}, \, a\neq b\) and let \(Z\) be the set of all families \(z\) indexed by \(I\) such that \(z_a \in X, z_b \in Y\). Now consider \(f : Z \leftrightarrow X \times Y\) such that \(f(z)=(z_a,z_b)\). This is the \textbf{\emph{generalized Cartesian product}}! There is further detail one can go into about families of sets and sets of all given families of sets, and even sets of those sets (including their Cartesian products), but these are so disgusting as to induce spontaneous retching in even the most tough-stomached mathematicians. So we will avoid it, if possible.
\section{Inverses and Composites}

Consider a function \(f : X \to  Y\), a natural mapping that might arise (when thinking on power sets) is some function \(F : \mcP (X) \to  \mcP(Y)\); just as each element of \(X\) is mapped to an element \(Y\) under \(f\), each subset of \(X\) maps to some subset of \(Y\) under \(F\), so if \(A \subset  X\) then \(\exists (A \mapsto \Img (f_A)) \forall  A \in  \mcP(X)\). Cool? I guess? What is to be done about this? In all honesty, I am not sure this factoid has much use other than being a pedagogical tool (though I am almost certainly wrong and I am sure someone will rush to correct me, wherein I might toil and languish over my incorrectness and stupidity, suffering immensely---but probably not, because I simply do not care), but for our cases we use it to introduce its behavior in the inverse.
\nt{Regarding notation---as we all too often are---\(\Img (f_A)\) is the same thing as \(f(A)\) when regarding images of subsets. I prefer the former notation as it feels a bit more natural to me.
}

\dfn{Inverses and inverse images}{Given a function \(f : X \to  Y\) we say \(f^{-1}\) is the \textbf{\emph{inverse}} of \(f\) given by \(f^{-1} : \mcP(Y) \to  \mcP(X)\) such that if \(B \subset  Y\), then
\[
    \Img(f^{-1}_B) = \{x \in  X \mid f(x) \in  B\},  \]
    is the \textbf{\emph{inverse image}} of \(B\) under \(f\)
}

This definition results in some immediate necessary and sufficient conditions: first, for \(f : X \to Y\) to exist, the inverse
image under \(f\) of each non-empty subset of \(Y\) must be a non-empty subset of \(X\); second, for \(f\) to be injective is
that the inverse image under \(f\) of each singleton in the range of \(f\) be a singleton in \(X\).

\nt{We often use \(f^{-1}\) for another purpose: for some \(f:X \to Y\), \(f^{-1} : \ran (f)\to X \implies f^{-1}(y)=x \iff f(x)=y\). This will be the most common use case for \(f^{-1}\) as the inverse function \(g : Y \to X\)
}


There are some connections between images and inverse images which I will quickly enumerate, but I will not provide proof for (for the proofs see Halmos,\textasciitilde{}pp. 59):
Let \(f\) be a function such that \(f : X \to A\) for each of the following,
\begin{enumerate}
\item \(B \subset  Y \implies f(f^{-1}(B))\subset B\);
\item \(f : X \surjto   Y, \, B \subset Y \implies f(f^{-1}(B))=B\);
\item \(A \subset X \implies A \subset  f^{-1}(f(A))\);
\item \(f : X \injto  Y , \, A \subset X  \implies f^{-1}(f(A))=A\);
\item If \(\{B_i\}\) is a family of subsets of \(Y\), then \(f^{-1}(\bigcup_{i \in  I} B_i ) =\bigcup_{i \in  I} f^{-1}(B_i)\) and \(f^{-1}(\bigcap_{i \in  I} B_i ) =\bigcap_{i \in  I} f^{-1}(B_i)\);
\item \(f^{-1}(Y-B)=X-f^{-1}(B)\) for each \(B \subset Y\).
\end{enumerate}

We now move onto \textbf{\emph{composites}}.
\dfn{Composition}{Let \(f : X \to Y\) and \(g : Y \to Z\) it is clear that we can make a map from \(X \to Z\) since the range of \(f\) can be the domain of \(g\). We can represent this new function---i.e., the \textbf{\emph{composite function}}---as either \(h :X \to Z\) or \(g \circ f\), and often simply just \(gf\). We read this right to left.
}

Notice that function composition is not always commutative, nor need it be. Also notice that this remains true under the special case that \(f :X \to Y\) and \(g : Y \to X\); functional composition, however, is always associative. There will be a strong proof for this later on (in Groups). Notice that if \(g : X\to Y\) for standard \(f\), if it is bijective we have an inverse function and identity map given by \(h\). Functional composition under inverses appears something like this: if \(f :X \to  Y\) and \(g: Y\to Z\), and \(f^{-1} : \mcP(Y) \to X\) and \(g^{-1} : \mcP(Z) \to \mcP(Y)\). Then \(g \circ f\) has an inverse given by \(f^{-1}\circ g^{-1}\). Notice this is not true for \(g^{-1}\circ f^{-1}\).
\subsection{Inverse/ converse and composite relations}
The main idea of this comes out of our exploration of inverse functions. Functions \emph{are} relations after all. Consider the \textbf{\emph{converse relation}} of \(R\) on \(X \times Y\) s.t. \(xRy\) and \(yR^{-1}x\). Similar things can be done for composites by generalizing composition to relations, observe: let \(S\) be a relation on \(Y \times Z\) and \(R\) be a relation on \(X \times Y\) such that \(xRy\) and \(ySz\). Now let \(T\) over \(X\times Z\) be given by \(S\circ R\) such that \(xTz\) exists \(\iff  \exists  y \in Y\) such that \(xRy\) and \(ySz\). Congrats! You made a \textbf\{\textit{composite relation}\}. If it is the case that \(R\) and \(S\) are functions s.t. \(R(x)=y, \, S(y)=z\) and \(S(R(x))=z \iff xTz\) which implies that functional composition is a special case of a greater abstraction called the \textbf\{\textit{relative product}\}(??).





\end{document}
